\documentclass[a4paper,12pt]{article}
\usepackage{amsmath, amssymb}
\usepackage{xcolor}
\usepackage{tcolorbox}
\usepackage{multicol}
\usepackage{geometry}
\graphicspath{ {../Images/} }
\hyphenpenalty=100000

\geometry{left=1.5cm, right=1.5cm, top=2cm, bottom=2cm}

\tcbuselibrary{theorems}

\definecolor{PurpleEx}{HTML}{5d478b}
\definecolor{GrayEx}{HTML}{f2f2f2}
\definecolor{GrayBG}{HTML}{d9d9d9}
\definecolor{BrownHist}{HTML}{8b5a2b}
\definecolor{BlackSav}{HTML}{000000}

\title{
\centering
\tcbset{colframe=BlackSav,colback=GrayBG}
\tcbox[tikznode]{
Chapitre 5 \\Dérivation et applications
}
}
\usepackage{tikz}
\author{}
\date{}

\newtcbtheorem
  []% init options
  {history}% name
  {Histoire}% title
  {%
    theorem name,
    colback=GrayBG,
    colframe=BrownHist,
    fonttitle=\bfseries,
  }% options
  {hist}% prefix

\newtcbtheorem
  []% init options
  {exercice}% name
  {Exercice}% title
  {%
    colback=GrayEx,
    colframe=PurpleEx,
    fonttitle=\bfseries,
  }% options
  {ex}% prefix


\begin{document}

\maketitle

\begin{multicols}{2}
\begin{exercice}[title=Exercice 1 $\quad \star \quad$ { [Calculer] }]{}{}
Dans chaque cas, déterminer le PGCD de $n$ et $m$.
\begin{enumerate}
    \item $n = 15$ et $m = 37$
    \item $n = 32$ et $m = 36$
    \item $n = 40$ et $m = 8$
\end{enumerate}
\end{exercice}

\begin{exercice}[title=Exercice 2 $\quad \star \quad$ { [Calculer] }]{}{}
Déterminer l’ensemble des entiers naturels $n \leq 100$ tels que $\text{PGCD}(n; 100) = 20$.
\end{exercice}

\begin{exercice}[title=Exercice 3 $\quad \star \quad$ { [Modéliser, Calculer] }]{}{}
Un fleuriste possède $105$ roses rouges et $63$ roses blanches. Il souhaite composer des bouquets tous identiques.  
En utilisant toutes les fleurs, quel est le plus grand nombre de bouquets qu’il pourra créer ?
\end{exercice}

\begin{exercice}[title=Exercice 4 $\quad \star \quad$ { [Modéliser, Calculer] }]{}{}
En astronomie, on observe deux étoiles $A$ et $B$ une même nuit. L’étoile $A$ apparaît périodiquement tous les $35$ jours tandis que l’étoile $B$ apparaît périodiquement tous les $77$ jours.  
Combien de jours faudra-t-il attendre pour pouvoir observer les étoiles $A$ et $B$ simultanément ?
\end{exercice}

\begin{exercice}[title=Exercice 5 $\quad \star \quad$ { [Modéliser, Calculer] }]{}{}
Une association organise une compétition sportive ; $144$ filles et $252$ garçons se sont inscrits. L’association désire répartir les inscrits en équipes mixtes. Le nombre de filles et de garçons doivent être les mêmes dans chaque équipe. Tous les inscrits doivent être dans une des équipes.  
\begin{enumerate}
    \item Quel est le nombre maximum d’équipes que cette association peut former ?
    \item Quelle est alors la composition de chaque équipe ?
\end{enumerate}
\end{exercice}

\begin{exercice}[title=Exercice 6 $\quad \star \star \star \quad$ { [Calculer, Chercher, Raisonner] }]{}{}
Soit $n$ un entier naturel impair. Déterminer $\text{PGCD}(n^2 + 2n - 3 ; n^2 + 4n + 3)$.
\end{exercice}

\begin{exercice}[title=Exercice 7 $\quad \star \quad$ { [Calculer] }]{}{}
Dans chaque cas, déterminer à l’aide de l’algorithme d’Euclide le PGCD de $n$ et $m$.
\begin{enumerate}
    \item $n = 150$ et $m = 357$
    \item $n = 72$ et $m = 448$
    \item $n = 300$ et $m = 2880$
    \item $n = 1027$ et $m = 1032$
\end{enumerate}
\end{exercice}

% Exercice 8
\begin{exercice}[title=Exercice 8 $\quad \star \quad$ { [Calculer] }]{}{}
    \begin{enumerate}
        \item Montrer que la fraction $\frac{1510}{503}$ est irréductible.
        \item Réduire la fraction $\frac{4389}{5544}$ en fraction irréductible.
    \end{enumerate}
    \end{exercice}
    
    % Exercice 9
    \begin{exercice}[title=Exercice 9 $\quad \star \quad$ { [Calculer] }]{}{}
    Dans chaque cas, déterminer une identité de Bézout entre $n$ et $m$.
    \begin{enumerate}
        \item $n = 17$ et $m = 5$
        \item $n = 25$ et $m = 72$
        \item $n = 36$ et $m = 300$
    \end{enumerate}
    \end{exercice}
    
    % Exercice 10
    \begin{exercice}[title=Exercice 10 $\quad \star \star \quad$ { [Calculer, Raisonner, Chercher] }]{}{}
    Soit $n \geq 1$ un entier naturel. Déterminer $\text{PGCD}(7n + 9 ; n + 1)$.
    \end{exercice}
    
    % Exercice 11
    \begin{exercice}[title=Exercice 11 $\quad \star \star \quad$ { [Calculer, Raisonner, Chercher] }]{}{}
    Montrer que pour tout entier naturel $n \geq 1$, la fraction $\frac{n}{n^2 + 1}$ est irréductible.
    \end{exercice}
    
    % Exercice 12
    \begin{exercice}[title=Exercice 12 $\quad \star \star \quad$ { [Représenter, Raisonner] }]{}{}
    Déterminer le PGCD de :
    \begin{enumerate}
        \item deux nombres entiers consécutifs ;
        \item deux nombres entiers pairs consécutifs ;
        \item deux nombres entiers impairs consécutifs.
    \end{enumerate}
    \end{exercice}
    
    % Exercice 13
    \begin{exercice}[title=Exercice 13 $\quad \star \star \quad$ { [Calculer, Chercher] }]{}{}
    Déterminer tous les couples $(x ; y)$ d'entiers naturels tels que $x + y = 42$ et $\text{PGCD}(x ; y) = 6$.
    \end{exercice}
    
    % Exercice 14
    \begin{exercice}[title=Exercice 14 $\quad \star \star \quad$ { [Calculer, Chercher] }]{}{}
    Déterminer tous les couples $(x ; y)$ d'entiers naturels tels que $xy = 100$ et $\text{PGCD}(x ; y) = 5$.
    \end{exercice}
    
    % Exercice 15
    \begin{exercice}[title=Exercice 15 $\quad \star \star \quad$ { [Calculer, Chercher] }]{}{}
    Déterminer tous les couples $(x ; y)$ d'entiers naturels tels que $xy = 439230$ et $\text{PGCD}(x ; y) = 121$.
    \end{exercice}
    
    % Exercice 16
    \begin{exercice}[title=Exercice 16 $\quad \star \star \star \quad$ { [Calculer, Chercher] }]{}{}
    Dans chaque cas, montrer que pour tout $n \in \mathbb{N}$, $u_n$ et $u_{n+1}$ sont premiers entre eux.
    \begin{enumerate}
        \item Suite de Fibonacci : $u_0 = 0$, $u_1 = 1$ et pour tout $n \geq 2$, $u_{n+2} = u_{n+1} + u_n$.
        \item Nombres de Mersenne : pour tout $n \geq 1$, $u_n = 2^n - 1$.
        \item Nombres de Fermat : pour tout $n \geq 1$, $u_n = 2^{2^n} + 1$.
    \end{enumerate}
    Indication : pour la question 3, on pourra commencer par calculer $u_{n+1} - (u_n - 2) \times u_n$.
    \end{exercice}
    
    % Exercice 17
    \begin{exercice}[title=Exercice 17 $\quad \star \star \quad$ { [Raisonner, Chercher] }]{}{}
    Montrer que pour tout entier naturel $n$ non nul, la fraction $\frac{n(2n+1)}{n+1}$ est irréductible.
    \end{exercice}
    
    % Exercice 18
    \begin{exercice}[title=Exercice 18 $\quad \star \star \star \quad$ { [Raisonner, Chercher] }]{}{}
    Soient $a$ et $b$ des entiers naturels non nuls. Montrer que si $\frac{a}{b}$ est une fraction irréductible, alors $\frac{ab}{a+b}$ est irréductible.
    \end{exercice}    
    
    % Exercice 19
\begin{exercice}[title=Exercice 19 $\quad \star\star\star \quad$ { [Raisonner, Chercher] }]{}{}
    Soit $p$ un nombre premier. Montrer que pour tout $k \in \{1, \dots, p-1\}$, $p$ divise $\binom{p}{k}$.
    \end{exercice}
    
    % Exercice 20
    \begin{exercice}[title=Exercice 20 $\quad \star \quad$ { [Calculer, Raisonner] }]{}{}
    Montrer que pour tout entier naturel $n$, $6$ divise $n(n+1)(n+2)$.
    \end{exercice}
    
    % Exercice 21
    \begin{exercice}[title=Exercice 21 $\quad \star \quad$ { [Calculer, Raisonner] }]{}{}
    Montrer que pour tout entier naturel $n$, $120$ divise $n(n^2-1)(n^2-4)$.
    \end{exercice}
    
    % Exercice 22
    \begin{exercice}[title=Exercice 22 $\quad \star\star \quad$ { [Raisonner] }]{}{}
    Montrer que $\sqrt{3}$ est un nombre irrationnel.
    \end{exercice}
    
    % Exercice 23
    \begin{exercice}[title=Exercice 23 $\quad \star\star\star \quad$ { [Chercher, Raisonner, Communiquer] }]{}{}
    Soit $n \in \mathbb{N}$. Montrer $\sqrt{n}$ est un entier naturel ou un irrationnel.
    \end{exercice}
    
    % Exercice 24
    \begin{exercice}[title=Exercice 24 $\quad \star \quad$ { [Calculer] }]{}{}
    Dans chaque cas, indiquer si $a$ est inversible modulo $n$ et, dans le cas échéant, déterminer son inverse.
    \begin{enumerate}
        \item $a = 15$ et $n = 31$
        \item $a = 84$ et $n = 291$
        \item $a = 286$ et $n = 315$
    \end{enumerate}
    \end{exercice}
    
    % Exercice 25
    \begin{exercice}[title=Exercice 25 $\quad \star\star \quad$ { [Calculer, Raisonner] }]{}{}
    Déterminer l'ensemble des entiers $n$ tels que $19n-1$ est divisible par $71$.
    \end{exercice}
    
    % Exercice 26
    \begin{exercice}[title=Exercice 26 $\quad \star \quad$ { [Calculer] }]{}{}
    Résoudre dans $\mathbb{Z}^2$ les équations suivantes :
    \begin{enumerate}
        \item $38x = 65y$
        \item $22x = 40y$
        \item $208x = 390y$
    \end{enumerate}
    \end{exercice}
    
    % Exercice 27
    \begin{exercice}[title=Exercice 27 $\quad \star \quad$ { [Calculer] }]{}{}
    Déterminer l'ensemble des entiers relatifs $n$ tels que :
    \begin{enumerate}
        \item $6n \equiv 0 \pmod{35}$
        \item $5n \equiv 10 \pmod{41}$
        \item $9n \equiv 0 \pmod{57}$
    \end{enumerate}
    \end{exercice}
    
    % Exercice 28
    \begin{exercice}[title=Exercice 28 $\quad \star\star \quad$ { [Calculer] }]{}{}
    Résoudre dans $\mathbb{Z}^2$ les équations suivantes :
    \begin{enumerate}
        \item $11x - 15y = 3$
        \item $24x + 33y = 7$
        \item $84x + 55y = 4$
    \end{enumerate}
    \end{exercice}
    
    % Exercice 29
    \begin{exercice}[title=Exercice 29 $\quad \star\star \quad$ { [Représenter, Raisonner] }]{}{}
    Dans un repère orthonormé $(O; \vec{i}, \vec{j}, \vec{k})$ de l'espace, on note :
    \begin{itemize}
        \item $P$ le plan passant par le point $A(5;3;7)$ et de vecteur normal $\vec{n}_1 = \begin{pmatrix} 4 \\ 1 \\ 3 \end{pmatrix}$ ;
        \item $Q$ le plan passant par le point $B(10;4;1)$ et de vecteur normal $\vec{n}_2 = \begin{pmatrix} 7 \\ 1 \\ 13 \end{pmatrix}$.
    \end{itemize}
    \begin{enumerate}
        \item Montrer que les plans $P$ et $Q$ ne sont pas parallèles.
        \item Existe-t-il des points à coordonnées entières situés sur la droite d'intersection des plans $P$ et $Q$ ? Si oui, préciser leurs coordonnées.
    \end{enumerate}
    \end{exercice}    

    % Exercice 30
\begin{exercice}[title=Exercice 30 $\quad \star \star \star \quad$ { [Calculer, Chercher, Raisonner] }]{}{}
    On se propose de résoudre le système de congruences suivant :
    \[
    (S) : 
    \begin{cases} 
    x \equiv 2 \pmod{15} \\
    x \equiv 8 \pmod{77}
    \end{cases}
    \]
    \begin{enumerate}
        \item L’objectif de cette question est de déterminer une des solutions de $(S)$.
        \begin{enumerate}
            \item Écrire une relation de Bézout entre $15$ et $77$.
            \item En déduire la valeur de deux entiers $a$ et $b$ tels que :
            \[
            \hspace{-4em}\begin{cases}
            a &\equiv\ 1\ [15] \\
            a &\equiv\ 0\ [77]
            \end{cases}
            \quad \text{et} \quad 
            \begin{cases}
            b &\equiv\ 0\ [15] \\
            b &\equiv\ 1\ [77].
            \end{cases}
            \]
            \item En déduire enfin la valeur d’une solution de $(S)$.
        \end{enumerate}
        \item Supposons que $x_1$ et $x_2$ soient deux solutions de $(S)$. Montrer qu’alors $x_1 \equiv x_2 \pmod{15 \times 77}$.
        \item Conclure en déterminant l’ensemble des solutions de $(S)$.
    \end{enumerate}
    \end{exercice}
    
    % Exercice 31
    \begin{exercice}[title=Exercice 31 $\quad \star \star \quad$ { [Calculer, Raisonner] }]{}{}
    En utilisant la même méthode que l’exercice 30, résoudre le système de congruences suivant :
    \[
    (S) : 
    \begin{cases} 
    x \equiv 17 \pmod{101} \\
    x \equiv 10 \pmod{35}
    \end{cases}
    \]
    \end{exercice}
    
    % Histoire : Problème des restes chinois
    \begin{history}[title=Histoire : Problème des restes chinois]{}
    La résolution d’un système de congruences est connu sous le nom de « problème des restes chinois ». La forme originale du problème apparaît dans le livre de Sun Zi (III\textsuperscript{e} siècle). Il est repris au XIII\textsuperscript{e} siècle par le mathématicien chinois Qin Jiushao dans son \emph{Traité mathématique en neuf chapitres} sous la forme suivante :
    \begin{quote}
    « Soient des objets en nombre inconnu. Si on les range par $3$ il en reste $2$. Si on les range par $5$, il en reste $3$ et si on les range par $7$, il en reste $2$. Combien a-t-on d’objets ? »
    \end{quote}
    En arithmétique, la résolution de système de congruences est en fait très courante étant donné qu’elle permet, en décomposant un entier $n$ en produit de facteurs premiers, de ramener l’étude de congruences modulo $n$ à des congruences modulo des nombres premiers (ou puissances de nombres premiers).
    \end{history}
    
    % Exercice 32
    \begin{exercice}[title=Exercice 32 $\quad \star \star \star \quad$ { [Modéliser, Calculer, Raisonner] }]{}{}
    En astronomie, on observe une étoile $A$ dans le ciel, puis, six jours plus tard, on observe une étoile $B$. Sachant que l’étoile $A$ apparaît périodiquement tous les $105$ jours et que l’étoile $B$ apparaît périodiquement tous les $81$ jours, combien de jours faudra-t-il attendre pour pouvoir observer les étoiles $A$ et $B$ simultanément ?
    \end{exercice}
    
    % Exercice 33
    \begin{exercice}[title=Exercice 33 $\quad \star \star \star \quad$ { [Modéliser, Calculer, Raisonner] }]{}{}
    Un bataillon de moins de $1000$ soldats est en manœuvre sur un champ d’exercices. S’il se met en colonnes de $8$, de $15$, ou de $25$, le dernier rang ne comporte dans les trois cas que $5$ soldats. Si le bataillon se mettait en colonnes de $11$, tous les rangs seraient-ils complets ?
    \end{exercice}
    
    % Exercice 34
    \begin{exercice}[title=Exercice 34 $\quad \star \star \star \quad$ { [Représenter] }]{}{}
    \begin{enumerate}
        \item Écrire une fonction algorithmique en langage Python permettant de calculer le PGCD de deux entiers $n$ et $m$ à l’aide de l’algorithme d’Euclide.
        \item Écrire ensuite une fonction algorithmique en langage Python permettant de déterminer une relation de Bézout entre deux entiers $n$ et $m$.
    \end{enumerate}
    \end{exercice}
    
\end{multicols}

\end{document}
